% Fill in \reportauthor before submission.
% Keep the first-person verification statements only if you personally performed those checks.
\documentclass[11pt]{article}
\usepackage[margin=1in]{geometry}
\usepackage{amsmath,amssymb}
\usepackage{graphicx}
\usepackage{booktabs}
\usepackage[section]{placeins}
\usepackage{microtype}
\usepackage[hidelinks]{hyperref}

\newcommand{\reportauthor}{Ning Zhang}
\newsavebox{\resultbox}
\newcommand{\resulttable}[3]{%
  \begin{table}[htbp]\centering
  \caption{#2}\label{#3}
  \begingroup\small\setlength{\tabcolsep}{4pt}
  \catcode`\^=12
  \sbox{\resultbox}{\input{#1}}
  \ifdim\wd\resultbox>\linewidth
    \resizebox{\linewidth}{!}{\usebox{\resultbox}}
  \else
    \usebox{\resultbox}
  \fi
  \endgroup
  \end{table}}
\newcommand{\resultfigure}[4][0.90]{%
  \begin{figure}[htbp]\centering
  \includegraphics[width=#1\textwidth]{#2}
  \caption{#3}\label{#4}
  \end{figure}}

\title{Homework 1: Discrete State Dynamic Programming}
\author{\reportauthor}
\date{ECONS 509, Fall 2026}

\begin{document}
\maketitle

\section{Overview}

I solve the two-state Huggett household problem at the fixed prices implied by
$K^0=5$. The final numerical specification uses modified Howard improvement,
sparse CSR power iteration, an exp-log asset grid on $[0,5]$, and $N=5000$.
The Bellman formulation and Euler derivation are given in
\texttt{manual/derivation.pdf}; generated tables and figures are included from
\texttt{results/}.

\section{Solver Comparison: Part (a)}

The baseline uses a uniform $[0,20]$ grid with $N=1000$ and zero initial values.
Convergence requires
\[
 d(x)=\max_m\frac{|T(x)_m-x_m|}{|x_m|+1}\leq10^{-8}.
\]
Pass counts include the terminating check, and solver times exclude the common
utility-matrix construction.

\resulttable{results/part_a_solvers.tex}
{Baseline solver comparison, $N=1000$ on $[0,20]$.}
{tab:solvers}

VFI, Howard, and modified Howard converge to identical policies in 379, 15,
and 19 passes. Modified Howard is the fastest converged method at 0.055136
seconds, so I select it. The contraction benchmark
$\log(10^{-8})/\log(0.96)=451.244$ is larger than the observed VFI count because
it is a contraction-based error-reduction benchmark, not the exact stopping
count for the relative Bellman metric.

\resulttable{results/part_a_theory_diagnostics.tex}
{Contraction and conditioning diagnostics, $J=I-\beta Q_G$.}
{tab:conditioning}

The fixed gradient step $\eta=0.0005$ was preregistered as a conservative choice:
at the baseline size, the bound on $\|I-\beta Q\|_2^2$ is below 2000, so the step
uses the corresponding inverse scale. Fixed-step gradient descent and Adam both
hit the 20,000-pass cap and remain unconverged. The requested benchmark
$(1-\beta)^{-2}=625$ is reported separately from the measured Euclidean
condition numbers. In this problem, the measured $\operatorname{cond}_2(J^TJ)$
is much larger, so 625 should not be read as the measured 2-norm condition
number of this nonnormal Jacobian.

\section{Invariant Distribution: Parts (b)--(d)}

Using the converged Howard policy, I compare power iteration, the eigenvector
method, and direct solution of $Q_G^T\pi=\pi$ with both dense and CSR
representations.

\resulttable{results/parts_bc_distributions.tex}
{Invariant-distribution methods for the common Howard policy.}
{tab:distributions}
\resulttable{results/parts_bc_agreement.tex}
{Sup-norm agreement among invariant distributions.}
{tab:agreement}

All six solutions pass the registered checks, with maximum disagreement about
$5.49\times10^{-12}$. Sparse power iteration is the fastest measured sparse
method at 0.003668 seconds, so I select power iteration with CSR for the
remaining experiments.

\section{Asset Range: Part (e)}

At $N=1000$, I test uniform grids with $k_{\max}=1,2,5,10,20,40$.

\resulttable{results/part_e_ranges.tex}
{Uniform asset-range comparison at $N=1000$.}
{tab:ranges}

The ranges ending at 1 and 2 are rejected because they place positive mass at
the top node. I select $k_{\max}=5$, the smallest tested range with zero top-node
mass and an upper support endpoint well below the grid boundary.

\section{Grid Construction: Part (f)}

At $N=1000$ on $[0,5]$, I compare the uniform grid with
\[
 k_n=\exp\!\left[\frac{n-1}{N-1}\log(1+k_{\max})\right]-1.
\]

\resulttable{results/part_f_grids.tex}
{Uniform and exp-log grids, $N=1000$ on $[0,5]$.}
{tab:grids}

Both grids pass all applicable checks. I select the exp-log grid because it has
lower mean and distribution-weighted Euler errors, although its unweighted
maximum residual is slightly higher.

\resultfigure[0.82]{results/stage4/f/part_f_euler_comparison.pdf}
{Euler-residual comparison for the uniform and exp-log grids.}
{fig:grid-euler}

\section{Grid Size: Part (g)}

I compare $N=100,500,1000,2000,5000$ under the selected configuration.

\resulttable{results/part_g_nodes.tex}
{Grid-size comparison on the exp-log $[0,5]$ grid.}
{tab:nodes}

All five solutions converge in 19 passes. Euler errors fall steadily as $N$
increases, while the $N=5000$ run remains computationally modest at about
1.95 seconds. I therefore select $N=5000$ because it gives the lowest tested
Euler errors at an acceptable runtime.

\resultfigure[0.72]{results/stage4/g/part_g_policies.pdf}
{Asset policies across grid sizes.}
{fig:policies}
\resultfigure[0.72]{results/stage4/g/part_g_joint_masses.pdf}
{Joint invariant probability masses across grid sizes.}
{fig:masses}
\FloatBarrier

\section{Euler Accuracy: Part (h)}

For slack borrowing states, the consumption-equivalent residual is
\[
 c^{EE}_{n,s}=\left[\beta(1+r^0)\sum_{s'}P_{ss'}
 (c^{\mathrm{next}}_{n,s,s'})^{-\sigma}\right]^{-1/\sigma},
 \qquad
 E_{n,s}=\left|1-\frac{c^{EE}_{n,s}}{c_{n,s}}\right|.
\]
For the slack-state set
\[
\mathcal A=\{(n,s):G(n,s)>0\},
\qquad
q_{n,s}=
\frac{\pi_{n,s}}
{\sum_{(i,j)\in\mathcal A}\pi_{i,j}}.
\]
The weighted maximum is
\[
\max_{(n,s)\in\mathcal A} q_{n,s}E_{n,s},
\]
while the supported slack maximum is
\[
\max_{\substack{(n,s)\in\mathcal A\\ \pi_{n,s}>10^{-12}}} E_{n,s}.
\]

\resulttable{results/part_h_accuracy.tex}
{Euler residuals at the selected $N=5000$ grid.}
{tab:accuracy}

The final mean, maximum, and conditional weighted mean residuals are
$2.23018\times10^{-4}$, $1.18456\times10^{-3}$, and
$1.55907\times10^{-4}$. No selected policy chooses the upper grid bound.

\resultfigure[0.82]{results/stage4/h/part_h_euler_residuals.pdf}
{Euler residuals at the selected grid size.}
{fig:final-errors}

Across the five saved grid sizes, the fitted log-log slopes are 1.013287 for the
mean error and 0.933748 for the maximum error.

\resulttable{results/part_h_scaling.tex}
{Log-log Euler-error scaling across five grids.}
{tab:scaling}

These slopes are close to one, which is consistent with approximately
first-order error reduction over the tested grid resolutions.

\resultfigure[0.82]{results/stage4/h/part_h_accuracy_scaling.pdf}
{Mean and maximum Euler errors against the maximum grid step.}
{fig:scaling}

\paragraph{Final checks.}
All nine selected-grid checks pass. A total of 202 check records pass,
including the 30 Stage-3 validation records, with two expected rejected trial
ranges. The reproduction check matches 330 saved arrays and 8,514 JSON numerical
values with zero numerical difference; the capped gradient and Adam runs remain
correctly labeled unconverged.

\section{Independent Verification}

\paragraph{Manual-kernel check.}
On October 1, 2026, I independently reran the manual-kernel comparison on the
$N=100$ uniform $[0,20]$ grid. Both the manual and implementation VFI runs
required 492 passes. The maximum value-function difference was
$3.71\times10^{-9}$ and there were zero policy mismatches, satisfying the
$10^{-8}$ comparison criterion.

\paragraph{Final Euler check.}
On October 1, 2026, I also independently checked the final $N=5000$ Euler
diagnostics. The mean and maximum residuals were 0.000223018 and 0.001184556,
matching the saved Stage-4 results to numerical precision.

These checks are separate from the execution agent's automated Stage-3 and
Stage-4 verification and are recorded in \texttt{log.md}.

\section{AI-Use Note}

I used Codex (GPT-6; exact build unavailable) through a desktop/local-folder
agent workflow to implement the production solvers and tests, run the numerical
experiments, and generate the results, tables, figures, and reproduction
evidence. I also used ChatGPT text chat to help draft and revise student-owned
workflow files, later manual-kernel maintenance edits, and this report draft. I
made the recorded runtime choices and reasons, reviewed the evidence, recorded
the verification in \texttt{log.md}, and am responsible for the final report and
Stage-5 commit. The report-check agent may correct only numbers, references, and
typographical errors.

\section{Conclusion}

Modified Howard with CSR power iteration on the $N=5000$ exp-log $[0,5]$ grid
passes the registered checks and yields small Euler residuals with approximately
first-order scaling over the tested resolutions.

\end{document}
